\section{هفته‌ی 8: مرور میان‌ترم: کل مسیر مدل‌سازی در یک مطالعه‌ی موردی}

\begin{quote}
\textbf{مرور فصل‌های 1 تا 3} · درس 1 ساعت (مطالعه‌ی موردی یکپارچه) · زمان
\lr{lab:} \textbf{آزمون میان‌ترم} (90 دقیقه) + مشاوره‌ی پروژه (30 دقیقه)
\end{quote}

\begin{center}\rule{0.5\linewidth}{0.5pt}\end{center}

\subsection{واژه‌های کلیدی}

واژه‌های مهم این درس با معادل فارسی‌شان (حتی اگر در درس‌های قبل آمده
باشند).

{\def\LTcaptype{none} % do not increment counter
\begin{longtable}[]{@{}
  >{\raggedright\arraybackslash}p{(\linewidth - 4\tabcolsep) * \real{0.3333}}
  >{\raggedright\arraybackslash}p{(\linewidth - 4\tabcolsep) * \real{0.3333}}
  >{\raggedright\arraybackslash}p{(\linewidth - 4\tabcolsep) * \real{0.3333}}@{}}
\toprule\noalign{}
\begin{minipage}[b]{\linewidth}\raggedright
\lr{English}
\end{minipage} & \begin{minipage}[b]{\linewidth}\raggedright
فارسی
\end{minipage} & \begin{minipage}[b]{\linewidth}\raggedright
توضیح کوتاه
\end{minipage} \\
\midrule\noalign{}
\endhead
\bottomrule\noalign{}
\endlastfoot
\lr{formal model} / \lr{formalize} & مدل فرمال / فرمال کردن & نوشتن
نظریه به شکل معادله یا الگوریتم \\
\lr{simulation} & شبیه‌سازی & تولید داده‌ی مصنوعی با اجرای مدل \\
\lr{fit} / \lr{fitting} & برازش & تنظیم پارامترها تا مدل با داده جور
شود \\
\lr{parameter recovery} & بازیابی پارامتر & آیا برازش، پارامترهای معلوم
شبیه‌سازی را برمی‌گرداند \\
\lr{likelihood} & درست‌نمایی & احتمال داده به‌عنوان تابعی از پارامترها \\
\lr{criterion (c)} & ملاک & آستانه‌ی تصمیم؛ نشان‌دهنده‌ی سوگیری پاسخ \\
\lr{least squares} & کمترین مربعات & روش برازش با کمینه کردن مجموع مربع
خطاها \\
\lr{cross-validation (CV)} & اعتبارسنجی متقابل \lr{(cross-validation)} &
برازش روی بخشی از داده و آزمون روی بخش کنارگذاشته \\
\lr{sensitivity (d′)} & حساسیت & فاصله‌ی توزیع علامت و نویز \\
\lr{parameter} & پارامتر & کمیت قابل تنظیم مدل که بین افراد یا شرایط فرق
می‌کند \\
\lr{SSE} / \lr{RMSD} & مجموع مربعات خطا / ریشه‌ی میانگین مربع انحراف &
معیارهای فاصله‌ی داده و پیش‌بینی \\
\lr{log-likelihood} / \lr{NLL} & لگاریتم درست‌نمایی / منفی آن & مقداری که
بهینه‌سازها کمینه یا بیشینه می‌کنند \\
\lr{estimate} / \lr{estimator} & برآورد / برآوردگر & مقدار به‌دست‌آمده
برای پارامتر / روش به دست آوردن آن \\
\lr{noise} & نویز & تغییرات تصادفی در شواهد یا در داده \\
\lr{deviance (G²)} & انحراف \lr{(deviance)} & معیار عدم برازش نسبت به
مدل اشباع‌شده \\
\lr{AIC} / \lr{BIC} & معیار اطلاعاتی آکائیکه / بیزی & برازش به‌علاوه‌ی
جریمه برای تعداد پارامتر \\
\lr{research design} & طرح آزمایش & شرایط، تعداد کوشش‌ها و نمونه‌ها \\
\lr{confusion matrix} & ماتریس درهم‌ریختگی \lr{(confusion matrix)} &
\lr{P(}مدل انتخاب‌شده \\
\lr{hit} & اصابت \lr{(hit)} & پاسخ «بله» وقتی علامت وجود دارد \\
\lr{response time (RT)} & زمان واکنش & فاصله‌ی ارائه‌ی محرک تا پاسخ \\
\lr{signal detection theory (SDT)} & نظریه‌ی تشخیص علامت & مدلی که حساسیت
را از سوگیری پاسخ جدا می‌کند \\
\lr{evidence} & شواهد & اطلاعاتی که برای تصمیم جمع می‌شود \\
\lr{exponential} / \lr{power function} & تابع نمایی / توانی & دو شکل
رایج برای منحنی فراموشی و یادگیری \\
\lr{training} / \lr{test data} & داده‌ی آموزش / آزمون & داده‌ی برازش /
داده‌ی کنارگذاشته برای سنجش \\
\end{longtable}
}

\begin{center}\rule{0.5\linewidth}{0.5pt}\end{center}

\subsection{اهداف یادگیری}

در پایان این هفته، دانشجویان می‌توانند:

\begin{enumerate}
\def\labelenumi{\arabic{enumi}.}
\tightlist
\item
  کل مسیر مدل‌سازی (فرمال‌سازی \ensuremath{\leftarrow} شبیه‌سازی \ensuremath{\leftarrow} برازش \ensuremath{\leftarrow} ارزیابی \ensuremath{\leftarrow} بازیابی \ensuremath{\leftarrow}
  تفسیر) را روی یک مسئله‌ی جدید بدون راهنمایی قدم‌به‌قدم \textbf{اجرا
  کنند}.
\item
  درست‌نمایی، روش برازش و روش ارزیابی مناسب برای یک مجموعه‌داده را
  \textbf{انتخاب کنند} و هر انتخاب را \textbf{توجیه کنند}.
\item
  مجموعه‌ای از برازش‌های مدل‌های رقیب، معیارهای اطلاعاتی و نتایج بازیابی را
  \textbf{تفسیر کنند} و یک نتیجه‌گیری قابل دفاع که محدودیت‌هایش را می‌گوید
  \textbf{بنویسند}.
\item
  فرمول‌های کلیدی هفته‌های 1 تا 7 (برگه‌ی فرمول میان‌ترم) را \textbf{به خاطر
  بیاورند} و \textbf{به کار ببرند}.
\end{enumerate}

\begin{center}\rule{0.5\linewidth}{0.5pt}\end{center}

\subsection{مقدمه}

در هفت هفته یک جعبه‌ابزار را تکه‌تکه ساختیم: مدل‌های فرمال (هفته‌های 1\lr{--}2)،
کمترین مربعات و بهینه‌سازی (هفته‌ی 3)، بیشینه درست‌نمایی (هفته‌ی 4)، عدم
قطعیت و چند شرکت‌کننده (هفته‌ی 5)، \lr{cross-validation} (هفته‌ی 6)،
معیارهای اطلاعاتی و بازیابی (هفته‌ی 7). هر \lr{lab} یکی دو تکه را تمرین
کرد. پژوهش واقعی همه‌ی آن‌ها را هم‌زمان، به ترتیب درست و با قضاوت درباره‌ی
اینکه چه چیزی مهم است، لازم دارد.

درس امروز کل مسیر را روی یک مسئله‌ی \textbf{جدید} اجرا می‌کند: آیا مطالعه‌ی
بیشتر یک کلمه آن را \emph{به‌یادماندنی‌تر} می‌کند (حساسیت بیشتر)، و اگر
بله، حساسیت با تکرارها چطور رشد می‌کند؟ این مطالعه‌ی موردی بهترین آمادگی
برای میان‌ترم هم هست که از شما می‌خواهد همین نوع تصمیم‌ها را بگیرید و توجیه
کنید.

\begin{center}\rule{0.5\linewidth}{0.5pt}\end{center}

\subsection{جزوه‌ی درس}

\paragraph{به زبان ساده}

این هفته همه‌ی تکه‌های نیمه‌ی اول درس را روی یک سؤال جدید کنار هم می‌گذارد:
\emph{آیا مطالعه‌ی بیشتر یک کلمه، بازشناسی بعدی آن را آسان‌تر می‌کند؟}
نظریه را به شکل یک مدل \lr{SDT} می‌نویسیم که در آن حساسیت \(d'\) به تعداد
تکرارها بستگی دارد. سه نسخه برای این وابستگی پیشنهاد می‌کنیم (بدون اثر،
افزایش هموار، هر الگویی)، هر سه را با بیشینه‌ی درست‌نمایی برازش می‌کنیم،
بررسی می‌کنیم که آیا هرکدام عدم برازش دارد (\(G^2\))، آن‌ها را مقایسه
می‌کنیم \lr{(LRT}، \lr{AIC}، \lr{BIC)}، بررسی می‌کنیم که روش در این طرح
آزمایش کار می‌کند (بازیابی)، و در پایان مدل برنده را به زبان ساده‌ی
روان‌شناختی تفسیر می‌کنیم، از جمله اینکه داده چه چیزهایی را
\emph{نمی‌تواند} به ما بگوید. اگر بتوانید این یک مثال را از اول تا آخر
دنبال کنید، نیمه‌ی اول درس را فهمیده‌اید.

\textbf{زمان‌بندی پیشنهادی (60 دقیقه):} نقشه‌ی مفهومی 10 · مطالعه‌ی موردی
40 · توضیح آزمون 10.

\subsubsection{1. نقشه‌ی مفهومی هفته‌های 1 تا 7}

{\def\LTcaptype{none} % do not increment counter
\begin{longtable}[]{@{}
  >{\raggedright\arraybackslash}p{(\linewidth - 6\tabcolsep) * \real{0.2500}}
  >{\raggedright\arraybackslash}p{(\linewidth - 6\tabcolsep) * \real{0.2500}}
  >{\raggedright\arraybackslash}p{(\linewidth - 6\tabcolsep) * \real{0.2500}}
  >{\raggedright\arraybackslash}p{(\linewidth - 6\tabcolsep) * \real{0.2500}}@{}}
\toprule\noalign{}
\begin{minipage}[b]{\linewidth}\raggedright
مرحله
\end{minipage} & \begin{minipage}[b]{\linewidth}\raggedright
سؤال
\end{minipage} & \begin{minipage}[b]{\linewidth}\raggedright
ابزارها
\end{minipage} & \begin{minipage}[b]{\linewidth}\raggedright
هفته
\end{minipage} \\
\midrule\noalign{}
\endhead
\bottomrule\noalign{}
\endlastfoot
فرمال‌سازی & نظریه دقیقاً چه ادعایی دارد؟ & پنج جزء یک مدل؛ پارامترها
به‌عنوان سازه‌های روان‌شناختی & 1\lr{--}2 \\
شبیه‌سازی & مدل چه کاری می‌تواند (و نمی‌تواند) بکند؟ & شبیه‌سازی مستقیم،
بررسی‌های تحلیلی، نمودار پیش‌بینی‌ها & 1\lr{--}2 \\
برازش & کدام مقادیر پارامتر داده‌ها را بهتر توضیح می‌دهند؟ & \lr{SSE} /
کمترین مربعات؛ \lr{NLL} / \lr{MLE}؛ شبکه، \lr{Nelder\lr{--}Mead}،
\lr{L-BFGS-B}؛ \lr{multi-start}؛ کران‌ها و \lr{reparameterization} &
3\lr{--}4 \\
سنجش عدم قطعیت & برآوردها چقدر دقیق‌اند؟ & \lr{SE} از \lr{Hessian
(Wald)}، \lr{profile likelihood}، \lr{bootstrap} & 5 \\
ارزیابی (مطلق) & آیا عدم برازش مدل از نویز بیشتر است؟ & باقی‌مانده‌ها،
\(G^2\) & 2\lr{--}4 \\
ارزیابی (نسبی) & کدام مدل بهتر تعمیم پیدا می‌کند؟ & \lr{CV}، \lr{LRT}،
\lr{AIC/AICc/BIC}، وزن‌ها & 6\lr{--}7 \\
بازیابی & آیا روش با این طرح آزمایش می‌تواند کار کند؟ & بازیابی پارامتر،
بازیابی مدل (ماتریس‌های درهم‌ریختگی و معکوس) & 7 \\
تفسیر & از نظر روان‌شناختی چه معنایی دارد و چه چیزی را نمی‌توانیم نتیجه
بگیریم؟ & معنای پارامترها، محدودیت‌های طرح آزمایش، خطاهای میانگین‌گیری &
همه \\
\end{longtable}
}

\textbf{راهنمای انتخاب درست‌نمایی:}

{\def\LTcaptype{none} % do not increment counter
\begin{longtable}[]{@{}
  >{\raggedright\arraybackslash}p{(\linewidth - 4\tabcolsep) * \real{0.3333}}
  >{\raggedright\arraybackslash}p{(\linewidth - 4\tabcolsep) * \real{0.3333}}
  >{\raggedright\arraybackslash}p{(\linewidth - 4\tabcolsep) * \real{0.3333}}@{}}
\toprule\noalign{}
\begin{minipage}[b]{\linewidth}\raggedright
داده
\end{minipage} & \begin{minipage}[b]{\linewidth}\raggedright
مدل نویز طبیعی
\end{minipage} & \begin{minipage}[b]{\linewidth}\raggedright
هفته
\end{minipage} \\
\midrule\noalign{}
\endhead
\bottomrule\noalign{}
\endlastfoot
شمارش از \(n\) (یادآوری، \lr{hit}، \lr{false alarm)} & دوجمله‌ای & 2،
4 \\
انتخاب‌های دودویی تکی & برنولی & 4 \\
پیوسته، تقریباً متقارن، با واریانس برابر (مثلاً میانگین زمان واکنش بلوک‌ها)
& گاوسی (= کمترین مربعات) & 4\lr{--}5 \\
\end{longtable}
}

\subsubsection{2. مطالعه‌ی موردی: آیا تکرار حساسیت حافظه را بالا می‌برد؟}

\textbf{طرح آزمایش.} در یک آزمایش حافظه‌ی بازشناسی، کلمه‌ها 1، 2 یا 4 بار
مطالعه می‌شوند. در آزمون، برای هر شرط تکرار 100 کلمه‌ی قدیمی و 100 کلمه‌ی
جدید (فریبنده یا \lr{lure)} ارائه می‌شود و شرکت‌کننده می‌گوید «قدیمی» یا
«جدید». چون کلمه‌های فریبنده هرگز مطالعه نشده‌اند، \textbf{یک} نرخ
\lr{false alarm} مشترک برای همه‌ی شرایط و \textbf{سه} نرخ \lr{hit} داریم.

\textbf{فرمال‌سازی \lr{(SDT} با ملاک مطلق).} کلمه‌های جدید شواهد
\(X \sim \mathcal{N}(0, 1)\) و کلمه‌های قدیمی‌ای که \(r\) بار مطالعه
شده‌اند \(X \sim \mathcal{N}(d'_r, 1)\) تولید می‌کنند؛ شرکت‌کننده اگر
\(X > \kappa\) باشد می‌گوید «قدیمی»، با یک ملاک \(\kappa\) برای همه‌ی
آیتم‌های آزمون (آیتم‌ها در هم آمیخته‌اند، پس شرکت‌کننده نمی‌داند هر آیتم از
کدام شرط است). پس

\[
H_r = \Phi(d'_r - \kappa), \qquad F = \Phi(-\kappa).
\]

\textbf{سه مدل رقیب برای وابستگی \(d'\) به تکرارها:}

{\def\LTcaptype{none} % do not increment counter
\begin{longtable}[]{@{}
  >{\raggedright\arraybackslash}p{(\linewidth - 6\tabcolsep) * \real{0.2500}}
  >{\raggedright\arraybackslash}p{(\linewidth - 6\tabcolsep) * \real{0.2500}}
  >{\raggedright\arraybackslash}p{(\linewidth - 6\tabcolsep) * \real{0.2500}}
  >{\raggedright\arraybackslash}p{(\linewidth - 6\tabcolsep) * \real{0.2500}}@{}}
\toprule\noalign{}
\begin{minipage}[b]{\linewidth}\raggedright
مدل
\end{minipage} & \begin{minipage}[b]{\linewidth}\raggedright
\(d'_r\)
\end{minipage} & \begin{minipage}[b]{\linewidth}\raggedright
\(k\)
\end{minipage} & \begin{minipage}[b]{\linewidth}\raggedright
ادعا
\end{minipage} \\
\midrule\noalign{}
\endhead
\bottomrule\noalign{}
\endlastfoot
\(\mathcal{M}_0\) (بی‌اثر) & \(\delta_0\) & 2 & تکرار حساسیت را تغییر
نمی‌دهد \\
\(\mathcal{M}_1\) (خطی\lr{--}لگاریتمی) & \(\delta_0 + \delta_1 \log_2 r\) & 3
& هر \emph{دو برابر شدن} مطالعه همان مقدار حساسیت اضافه می‌کند \\
\(\mathcal{M}_2\) (آزاد) & \(d'_1, d'_2, d'_4\) & 4 & هر الگویی
(اشباع‌شده: 4 پارامتر برای 4 نرخ) \\
\end{longtable}
}

\(\mathcal{M}_0 \subset \mathcal{M}_1 \subset \mathcal{M}_2\)، پس همه‌ی
جفت‌ها \lr{nested} هستند.

\textbf{درست‌نمایی.} دوجمله‌ای برای هر یک از چهار شمارش (سه شمارش
\lr{hit}، یک شمارش \lr{false alarm).}

کد \lr{lab} در ادامه همه‌ی مراحل را اجرا می‌کند. دانشجویان آن را برای
آمادگی آزمون در خانه اجرا می‌کنند؛ در درس خروجی‌ها مرور می‌شوند.

\paragraph{خواندن ریاضی}

\begin{itemize}
\tightlist
\item
  \(H_r = \Phi(d'_r - \kappa)\) و \(F = \Phi(-\kappa)\): همان \lr{SDT}
  هفته‌ی 2، ولی ملاک به‌صورت یک موقعیت \textbf{مطلق} \(\kappa\) روی محور
  شواهد نوشته شده (از مرکز منحنی «کلمه‌ی جدید» سنجیده می‌شود) به جای اینکه
  نسبت به نقطه‌ی وسط سنجیده شود. کلمه‌های جدید: منحنی در 0، پس \(F\) مساحت
  سمت راست \(\kappa\) است، یعنی \(\Phi(-\kappa)\). کلمه‌های قدیمی که
  \(r\) بار مطالعه شده‌اند: منحنی در \(d'_r\)، پس
  \(H_r = \Phi(d'_r - \kappa)\).
\item
  یک \(\kappa\) برای همه‌ی شرایط: همه‌ی کلمه‌های آزمون با هم مخلوط‌اند، پس
  شرکت‌کننده نمی‌تواند برای کلمه‌هایی که یک بار و کلمه‌هایی که چهار بار
  مطالعه شده‌اند آستانه‌ی متفاوتی به کار ببرد.
\item
  \(\log_2 r\): لگاریتم در پایه‌ی 2، پس \(\log_2 1 = 0\)،
  \(\log_2 2 = 1\) و \(\log_2 4 = 2\). در مدل \(\mathcal{M}_1\)،
  \(d'_r = \delta_0 + \delta_1\log_2 r\) یعنی: \(\delta_0\) حساسیت بعد
  از یک بار مطالعه است، و هر \emph{دو برابر} شدن مطالعه \(\delta_1\)
  اضافه می‌کند.
\item
  \textbf{عددهای مثال حل‌شده:} \(\hat\kappa = 1.04\) می‌دهد
  \(F = \Phi(-1.04) = 0.15\) \ensuremath{\checkmark}. با \(\hat d'_1 = 0.82\):
  \(H_1 = \Phi(0.82 - 1.04) = \Phi(-0.22) = 0.41\)، نزدیک مقدار
  مشاهده‌شده‌ی 0.40 \ensuremath{\checkmark}.
\end{itemize}

\paragraph{برای یک مهندس}

مسیر نقشه‌ی مفهومی همان روند \textbf{شناسایی سیستم \lr{(system
identification)}} است که می‌شناسید: ساختارهای مدل نامزد را انتخاب کن،
پارامترها را تخمین بزن \lr{(ML)}، اعتبارسنجی کن (آزمون باقی‌مانده‌ها \ensuremath{\approx}
\(G^2\))، مرتبه‌ی مدل را انتخاب کن \lr{(LRT/AIC/BIC} \ensuremath{\approx}
\lr{GLRT/AIC/MDL)}، و روش را روی داده‌ی شبیه‌سازی‌شده با حقیقت معلوم وارسی
کن (بازیابی). دو تفاوت ارزش تأکید در کلاس را دارند: «سیستم» یک آدم است
که خود پارامترهایش نتیجه‌ی علمی‌اند، و مقدار داده با معیارهای مهندسی خیلی
کم است (400 پاسخ دودویی)، پس عدم قطعیت و بازیابی خیلی مهم‌ترند.

\paragraph{برای مدرس}

\begin{itemize}
\tightlist
\item
  \textbf{جدول نقشه‌ی مفهومی را ستون اصلی جلسه کنید.} برای هر سطر، خروجی
  متناظر در مطالعه‌ی موردی را نشان دهید (مثلاً «ارزیابی مطلق» \ensuremath{\leftarrow} ستون
  \(G^2\)). دانشجوها می‌بینند هر ابزار جایی دارد.
\item
  \textbf{پیش از خودتان از کلاس بخواهید تفسیر کند.} جدول نتایج را بدون
  پاراگراف تفسیر نشان دهید و بگذارید گروه‌های کوچک دو جمله بنویسند: «چه
  نتیجه‌ای می‌گیریم و چه نتیجه‌ای نمی‌توانیم بگیریم؟» بعد با متن مقایسه
  کنید. این دقیقاً مهارتی است که میان‌ترم می‌سنجد.
\item
  \textbf{توضیح امتحان:} به دانشجوها بگویید امتحان به \emph{تصمیم‌های
  مستدل} (کدام درست‌نمایی، کدام معیار، چرا بازیابی مهم است) بیشتر از
  محاسبه نمره می‌دهد. یک سؤال به سبک امتحان از بخش خودآزمایی نشان دهید.
\item
  \textbf{مشاوره‌ی پروژه (وقت \lr{lab):}} از هر تیم سه سؤال بپرسید: سؤال
  شناختی شما چیست؟ کدام پارامتر یا مقایسه‌ی مدل به آن جواب می‌دهد؟ آیا طرح
  آزمایش شما می‌تواند آن را بازیابی کند؟ بیشتر پیشنهادهای ضعیف در سؤال
  سوم شکست می‌خورند.
\end{itemize}

\begin{center}\rule{0.5\linewidth}{0.5pt}\end{center}

\subsection{مثال حل‌شده}

نتایج مطالعه‌ی موردی برای یک شرکت‌کننده‌ی شبیه‌سازی‌شده (مدل واقعی
\(\mathcal{M}_1\) با \(\delta_0 = 0.8\)، \(\delta_1 = 0.6\)،
\(\kappa = 1.0\))، از کد بخش \lr{lab.} اعداد گرد شده‌اند.

\textbf{قدم 1: داده.}

{\def\LTcaptype{none} % do not increment counter
\begin{longtable}[]{@{}llll@{}}
\toprule\noalign{}
شرط & کلمه‌های قدیمی & پاسخ «قدیمی» & نرخ \\
\midrule\noalign{}
\endhead
\bottomrule\noalign{}
\endlastfoot
\(r = 1\) & 100 & 40 & 0.40 \\
\(r = 2\) & 100 & 66 & 0.66 \\
\(r = 4\) & 100 & 80 & 0.80 \\
جدید (فریبنده) & 100 & 15 & 0.15 \\
\end{longtable}
}

\textbf{قدم 2: برازش و ارزیابی} \lr{(MLE}، \(N = 400\) پاسخ):

{\def\LTcaptype{none} % do not increment counter
\begin{landscape}\small
\begin{longtable}[]{@{}
  >{\raggedright\arraybackslash}p{(\linewidth - 18\tabcolsep) * \real{0.1000}}
  >{\raggedright\arraybackslash}p{(\linewidth - 18\tabcolsep) * \real{0.1000}}
  >{\raggedright\arraybackslash}p{(\linewidth - 18\tabcolsep) * \real{0.1000}}
  >{\raggedright\arraybackslash}p{(\linewidth - 18\tabcolsep) * \real{0.1000}}
  >{\raggedright\arraybackslash}p{(\linewidth - 18\tabcolsep) * \real{0.1000}}
  >{\raggedright\arraybackslash}p{(\linewidth - 18\tabcolsep) * \real{0.1000}}
  >{\raggedright\arraybackslash}p{(\linewidth - 18\tabcolsep) * \real{0.1000}}
  >{\raggedright\arraybackslash}p{(\linewidth - 18\tabcolsep) * \real{0.1000}}
  >{\raggedright\arraybackslash}p{(\linewidth - 18\tabcolsep) * \real{0.1000}}
  >{\raggedright\arraybackslash}p{(\linewidth - 18\tabcolsep) * \real{0.1000}}@{}}
\toprule\noalign{}
\begin{minipage}[b]{\linewidth}\raggedright
مدل
\end{minipage} & \begin{minipage}[b]{\linewidth}\raggedright
\(k\)
\end{minipage} & \begin{minipage}[b]{\linewidth}\raggedright
\lr{NLL}
\end{minipage} & \begin{minipage}[b]{\linewidth}\raggedright
\(G^2\) \lr{(df}، \(p\))
\end{minipage} & \begin{minipage}[b]{\linewidth}\raggedright
\lr{AIC}
\end{minipage} & \begin{minipage}[b]{\linewidth}\raggedright
\lr{BIC}
\end{minipage} & \begin{minipage}[b]{\linewidth}\raggedright
\(w_{\text{\lr{AIC}}}\)
\end{minipage} & \begin{minipage}[b]{\linewidth}\raggedright
\(w_{\text{\lr{BIC}}}\)
\end{minipage} & \begin{minipage}[b]{\linewidth}\raggedright
\(\hat d'_{1,2,4}\)
\end{minipage} & \begin{minipage}[b]{\linewidth}\raggedright
\(\hat\kappa\)
\end{minipage} \\
\midrule\noalign{}
\endhead
\bottomrule\noalign{}
\endlastfoot
\(\mathcal{M}_0\) بی‌اثر & 2 & 27.27 & 35.55 (2، \lr{\textless} .001) &
58.54 & 66.52 & 0.00 & 0.00 & 1.34، 1.34، 1.34 & 1.04 \\
\(\mathcal{M}_1\) خطی\lr{--}لگاریتمی & 3 & 9.77 & 0.54 (1، .46) &
\textbf{25.53} & \textbf{37.51} & \textbf{0.67} & \textbf{0.94} & 0.82،
1.37، 1.93 & 1.04 \\
\(\mathcal{M}_2\) آزاد & 4 & 9.50 & 0 (اشباع‌شده) & 26.99 & 42.96 & 0.33
& 0.06 & 0.78، 1.45، 1.88 & 1.04 \\
\end{longtable}
\end{landscape}
}

آزمون‌های نسبت درست‌نمایی: \(\mathcal{M}_0\) در برابر \(\mathcal{M}_1\):
\(\Lambda = 35.01\)، \lr{df} = 1، \(p < .001\)؛ \(\mathcal{M}_1\) در
برابر \(\mathcal{M}_2\): \(\Lambda = 0.54\)، \lr{df} = 1، \(p = .46\).

\textbf{بازیابی با این طرح آزمایش} (کد \lr{lab}، بلوک دوم): پارامترهای
\(\mathcal{M}_1\) خوب بازیابی می‌شوند (\(r = 0.89\) تا \(0.93\)) و
\lr{BIC} مدل تولیدکننده را در 87 تا 98 درصد شبیه‌سازی‌ها شناسایی می‌کند.

\textbf{قدم 3: تفسیر.} تکرار به‌وضوح حساسیت را بالا می‌برد: مدل بی‌اثر
به‌شدت بد برازش می‌شود (\(G^2 = 35.5\)، \(p < .001\)) و همه‌ی معیارها آن را
رد می‌کنند. مدل خطی\lr{--}لگاریتمی با یک پارامتر کمتر به خوبی مدل اشباع‌شده
برازش می‌شود \lr{(LRT} با \(p = .46\)) و هم \lr{AIC} و هم \lr{BIC} آن را
ترجیح می‌دهند. برآورد \(\hat\delta_1 \approx 0.55\) یعنی هر دو برابر شدن
مطالعه تقریباً نیم واحد \(d'\) اضافه می‌کند. نتایج بازیابی نشان می‌دهد که
با این طرح آزمایش، هم به پارامترها و هم به انتخاب مدل می‌شود اعتماد کرد.
هشدارها: این یک شرکت‌کننده‌ی شبیه‌سازی‌شده است؛ با فقط سه سطح تکرار، یک روند
خطی\lr{--}لگاریتمی را نمی‌شود از خیلی منحنی‌های صعودی دیگر (مثلاً یک تابع توانی
از \(r\)) جدا کرد؛ و یک ملاک واحد \(\kappa\) فرض می‌کند شرکت‌کننده پیش از
پاسخ نمی‌داند یک آیتم چند بار مطالعه شده است.

\begin{center}\rule{0.5\linewidth}{0.5pt}\end{center}

\subsection{\texorpdfstring{\lr{lab} (2 ساعت): آزمون میان‌ترم + مشاوره‌ی
پروژه}{lab (2 ساعت): آزمون میان‌ترم + مشاوره‌ی پروژه}}

\subsubsection{\texorpdfstring{آزمون میان‌ترم (90 دقیقه، در زمان
\lr{lab)}}{آزمون میان‌ترم (90 دقیقه، در زمان lab)}}

\begin{itemize}
\tightlist
\item
  \textbf{شکل:} کتبی، کتاب‌بسته. یک \textbf{برگه‌ی فرمول} (در ادامه) داده
  می‌شود. بدون لپ‌تاپ: سؤال‌ها شامل خواندن تکه‌کدهای کوتاه و تفسیر
  خروجی‌هاست، نه نوشتن کد از صفر.
\item
  \textbf{محدوده:} هفته‌های 1 تا 7 (فصل‌های 1 تا 3)، با تأکید بر \emph{فهم
  و تفسیر} (هدف ویژه‌ی 2 درس).
\item
  \textbf{انواع سؤال:} (1) سؤال‌های مفهومی کوتاه؛ (2) یک محاسبه‌ی دستی
  (مثلاً \(d'\) و \(c\)، یک درست‌نمایی یا \lr{NLL}، \lr{AIC/BIC} و وزن‌ها)؛
  (3) خواندن کد و پیش‌بینی یا توضیح خروجی آن؛ (4) تفسیر یک جدول برازش،
  نمودار باقی‌مانده‌ها یا ماتریس بازیابی/درهم‌ریختگی؛ (5) یک سؤال کوتاه
  طراحی («چه چیزی را تغییر می‌دادید و چرا؟»).
\item
  \textbf{وزن:} 3 نمره از 20.
\item
  خود آزمون میان‌ترم با پاسخ‌های نمونه در مواد ارزشیابی درس آمده است.
\end{itemize}

\textbf{برگه‌ی فرمولی که در آزمون داده می‌شود:}

\[
\begin{aligned}
&d' = \Phi^{-1}(H) - \Phi^{-1}(F), \quad c = -\tfrac12[\Phi^{-1}(H) + \Phi^{-1}(F)], \quad H = \Phi(d'/2 - c), \quad F = \Phi(-d'/2 - c)\\
&f_{\exp}(t) = A e^{-\lambda t}, \quad f_{\text{\lr{pow}}}(t) = A(1+t)^{-d}, \quad P_c^{\text{\lr{RW}}} = \frac{1-r^z}{1-r^a},\ r = \frac{1-v}{1+v}, \quad \mathbb{E}[T] = \frac{aP_c - z}{v}\\
&\mathrm{SSE} = \textstyle\sum_i (y_i - \hat y_i)^2, \quad \mathrm{RMSD} = \sqrt{\mathrm{SSE}/N}\\
&\ell(\boldsymbol\theta) = \textstyle\sum_i \big[\ln\binom{n_i}{y_i} + y_i\ln f_i + (n_i - y_i)\ln(1-f_i)\big], \quad \mathrm{NLL} = -\ell\\
&G^2 = 2(\ell_{\text{\lr{sat}}} - \hat\ell), \quad \Lambda = 2(\hat\ell_1 - \hat\ell_0) \sim \chi^2_{\Delta k}\\
&\mathrm{SE} \approx \sqrt{\mathrm{diag}(\mathbf{H}^{-1})}, \quad \text{\lr{Wald CI: }} \hat\theta \pm 1.96\,\mathrm{SE}, \quad \text{\lr{profile CI: }} 2(\mathrm{NLL}_{\text{\lr{prof}}} - \mathrm{NLL}_{\min}) \le 3.84\\
&\mathrm{AIC} = 2\,\mathrm{NLL} + 2k, \quad \mathrm{AICc} = \mathrm{AIC} + \tfrac{2k(k+1)}{N-k-1}, \quad \mathrm{BIC} = 2\,\mathrm{NLL} + k\ln N, \quad w_i = \tfrac{e^{-\Delta_i/2}}{\sum_j e^{-\Delta_j/2}}
\end{aligned}
\]

مقادیر بحرانی: \(\chi^2_{1,.95} = 3.84\)، \(\chi^2_{2,.95} = 5.99\)،
\(\chi^2_{3,.95} = 7.81\). \lr{quantile}ها:
\(\Phi^{-1}(0.84) \approx 1.0\)، \(\Phi^{-1}(0.975) = 1.96\). لگاریتم‌ها:
\(\ln 2 = 0.693\)، \(\ln 10 = 2.303\)، \(\ln 100 = 4.605\)،
\(\ln 400 = 5.991\).

\subsubsection{مشاوره‌ی پروژه (30 دقیقه، بعد از آزمون)}

هر گروه 5 تا 10 دقیقه با مدرس یا دستیار آموزشی درباره‌ی ایده‌ی پروژه‌اش پیش
از موعد پیشنهاد (ابتدای هفته‌ی 9) صحبت می‌کند. یک پاراگراف توصیف از سؤال،
مدل و داده‌تان بیاورید.

\subsubsection{مطالعه‌ی شخصی: مطالعه‌ی موردی در کد (پیش از آزمون اجرا
کنید)}

\begin{quote}
\textbf{\lr{Code} · کد:} \lr{\nolinkurl{code/week08.py}}, \lr{cell} 1 ·
سلول 1
\end{quote}

\begin{quote}
\textbf{\lr{Code} · کد:} \lr{\nolinkurl{code/week08.py}}, \lr{cell} 2 ·
سلول 2
\end{quote}

\textbf{خروجی مورد انتظار (بلوک اول):}

\Needspace{13\baselineskip}
\begin{latin}
\begin{Verbatim}
hits per condition: [40 66 80]  false alarms: 15
           k     NLL      G2  df  p(G2)     AIC     BIC            d'_1,2,4  kappa  w_AIC  w_BIC
model
M0_none    2  27.269  35.549   2  0.000  58.538  66.521  [1.34, 1.34, 1.34]   1.04  0.000  0.000
M1_loglin  3   9.766   0.543   1  0.461  25.533  37.507  [0.82, 1.37, 1.93]   1.04  0.674  0.938
M2_free    4   9.495   0.000   0    NaN  26.990  42.956  [0.78, 1.45, 1.88]   1.04  0.326  0.062

LRT M0 vs M1: 35.01 (df 1, p = 0.0000);  M1 vs M2: 0.54 (df 1, p = 0.461)
\end{Verbatim}
\end{latin}

\textbf{خروجی مورد انتظار (بلوک دوم، حدود 1 دقیقه):}

\Needspace{14\baselineskip}
\begin{latin}
\begin{Verbatim}
M1 recovery delta0 : r = 0.895, RMSE = 0.181
M1 recovery delta1 : r = 0.930, RMSE = 0.113
M1 recovery kappa  : r = 0.887, RMSE = 0.142

BIC model recovery, P(selected | generated):
           M0_none  M1_loglin  M2_free
M0_none       0.98       0.02     0.00
M1_loglin     0.05       0.93     0.02
M2_free       0.10       0.03     0.87
\end{Verbatim}
\end{latin}

یادداشت درباره‌ی مولد \(\mathcal{M}_2\) در بازیابی مدل: سه مقدار \(d'\)
آن طوری کشیده می‌شوند که \(d'_4\) \emph{بین} \(d'_1\) و \(d'_2\) قرار
بگیرد؛ یک الگوی غیریکنوا که \(\mathcal{M}_1\) نمی‌تواند تولید کند. در غیر
این صورت داده‌های \(\mathcal{M}_2\) اغلب کاملاً با \(\mathcal{M}_1\)
سازگار می‌بودند و «اشتباه گرفتن‌ها» بی‌معنی می‌شد.

\textbf{سؤال‌های خودسنجی درباره‌ی مطالعه‌ی موردی} (در درس بحث می‌شوند): 1.
چرا فقط یک نرخ \lr{false alarm} داریم، و این مدل \lr{SDT} را در مقایسه
با هفته‌ی 2 چطور تغییر می‌دهد؟ 2. هر معیار کدام مدل را ترجیح می‌دهد، و آیا
\lr{LRT}ها با آن موافق‌اند؟ 3. آیا \(\delta_1\) با این طرح آزمایش قابل
بازیابی است؟ آیا \(\mathcal{M}_1\) از \(\mathcal{M}_0\) و
\(\mathcal{M}_2\) قابل تشخیص است؟ 4. نتیجه‌ی روان‌شناختی را در یک جمله،
همراه با مهم‌ترین هشدارش، بنویسید.

\begin{center}\rule{0.5\linewidth}{0.5pt}\end{center}

\subsection{\texorpdfstring{ارتباط با هوش مصنوعی و
\lr{LLM}ها}{ارتباط با هوش مصنوعی و LLMها}}

\textbf{\lr{chatbot} به‌عنوان هم‌درس: اعتماد کن، ولی بررسی کن.} بسیاری از
دانشجوها از \lr{LLM}ها برای توضیح مفاهیم و بررسی جواب‌های امتحانی استفاده
می‌کنند. این می‌تواند کمک کند، ولی \lr{LLM}ها می‌توانند محاسبه‌هایی روان،
مطمئن و غلط تولید کنند. با جواب یک \lr{chatbot} مثل پیش‌بینی یک مدل رفتار
کنید: آن را با چیزی که خودتان می‌توانید حساب کنید بسنجید. کد پایین یک
جواب ازپیش‌آماده‌ی \lr{chatbot} به سؤال \lr{B3} میان‌ترم \lr{(AIC} برای سه
مدل) را بررسی می‌کند. \lr{chatbot} به جای \(2\,\mathrm{NLL} + 2k\) از
\(2\,\mathrm{NLL} + k\) استفاده کرده.

\begin{quote}
\textbf{\lr{Code} · کد:} \lr{\nolinkurl{code/week08.py}}, \lr{cell} 3 ·
سلول 3
\end{quote}

\textbf{خروجی مورد انتظار:}

\Needspace{9\baselineskip}
\begin{latin}
\begin{Verbatim}
M1: chatbot 222.0, correct 224.0  WRONG
M2: chatbot 218.0, correct 221.0  WRONG
M3: chatbot 218.8, correct 223.8  WRONG
Akaike weights, chatbot: [0.07 0.55 0.37]  correct: [0.15 0.68 0.17]
\end{Verbatim}
\end{latin}

\lr{chatbot} اتفاقاً برنده‌ی درست (\(\mathcal{M}_2\)) را انتخاب می‌کند، ولی
همه‌ی عددها غلط‌اند و پشتیبانی از پیچیده‌ترین مدل را دو برابر می‌کند. مصحح
نمره‌ی کمی می‌داد، و در یک تحلیل واقعی این خطا به نفع \lr{overfitting}
تمام می‌شد.

\textbf{امتحانش کنید (بدون نیاز به \lr{API).}} - \ensuremath{\star} کدام خطا در فرمول
دقیقاً عددهای \lr{chatbot} را تولید می‌کند؟ \lr{BIC} با همین خطا چه شکلی
می‌شد؟ - \ensuremath{\star}\ensuremath{\star} سه «مورد آزمون» (ورودی با جواب معلوم) کوچک برای فرمول‌های
\lr{SDT} هفته‌ی 2 بنویسید که بتوانید با آن هر جوابی را، چه از
\lr{chatbot} چه از هم‌کلاسی، بررسی کنید. - یادآوری: در خود امتحان هیچ
ابزار هوش مصنوعی مجاز نیست، و قانون درس درباره‌ی استفاده از هوش مصنوعی در
تکلیف‌ها در \lr{\nolinkurl{assessments/README.md}} آمده است.

\begin{center}\rule{0.5\linewidth}{0.5pt}\end{center}

\subsection{اشتباهات رایج و برداشت‌های نادرست (10 مورد اول از هفته‌های 1
تا 7)}

\begin{enumerate}
\def\labelenumi{\arabic{enumi}.}
\tightlist
\item
  برازش خوب را شاهد درستی مدل دانستن (هفته‌های 3، 6).
\item
  فراموش کردن علامت منفی: بهینه‌سازها \textbf{منفی} لگاریتم درست‌نمایی را
  کمینه می‌کنند (هفته‌ی 4).
\item
  محدود نکردن احتمال‌ها، که باعث می‌شود \(\ln 0\) برازش را از کار بیندازد
  (هفته‌ی 4).
\item
  اعتماد به \lr{\nolinkurl{success=True}}، یک شروع تنها، یا
  \lr{Nelder\lr{--}Mead} کران‌داری که روی یک کران فرو ریخته است (هفته‌های
  3\lr{--}4).
\item
  کمترین مربعات برای شمارش‌ها یا انتخاب‌ها بدون فکر کردن به مدل نویز
  (هفته‌های 3\lr{--}4).
\item
  گزارش برآوردها بدون عدم قطعیت، یا استفاده از بازه‌های \lr{Wald} نزدیک
  کران‌ها (هفته‌ی 5).
\item
  نتیجه‌گیری درباره‌ی افراد از منحنی‌های میانگین‌گیری‌شده (هفته‌ی 5).
\item
  نشت داده‌ی آزمون به آموزش در \lr{cross-validation} (هفته‌ی 6).
\item
  مقایسه‌ی مدل‌های \lr{non-nested} با \lr{LRT}، یا خواندن مقادیر \lr{AIC}
  به‌تنهایی به‌جای تفاوت‌هایشان (هفته‌ی 7).
\item
  تفسیر پارامترها یا مقایسه‌ی مدل‌ها بدون مطالعه‌ی بازیابی با طرح آزمایش
  واقعی (هفته‌ی 7).
\end{enumerate}

\begin{center}\rule{0.5\linewidth}{0.5pt}\end{center}

\subsection{خودآزمایی (سؤال‌های تمرینی به سبک آزمون)}

\textbf{1. (مفهومی)} در سه جمله توضیح دهید چرا مدل نمایی با مجانب
داده‌های هفته‌ی 2 را بهتر از مدل توانی برازش کرد، با اینکه مدل توانی داده
را تولید کرده بود. \textbf{پاسخ} مدل مجانبی یک پارامتر اضافه دارد، پس
می‌تواند خودش را با نویز خاص نمونه، اینجا نقاط صاف‌نمای فاصله‌های طولانی،
تطبیق دهد. برازش آموزش بهتر برای مدل‌های منعطف‌تر تضمین‌شده است و شاهد
درستی نیست. برای قضاوت به معیارهای تعمیم \lr{(CV}، \lr{AIC/BIC)} و
مطالعات بازیابی نیاز داریم.

\textbf{2. (محاسبه)} یک مشاهده‌گر \(H = 0.84\) و \(F = 0.16\) دارد.
\(d'\) و \(c\) را حساب کنید (از \(\Phi^{-1}(0.84) \approx 1.0\) استفاده
کنید). \textbf{پاسخ} \(d' = 1.0 - (-1.0) = 2.0\)؛
\(c = -\tfrac12(1.0 + (-1.0)) = 0\)، پس مشاهده‌گر بی‌سوگیری است.

\textbf{3. (محاسبه)} دو مدل روی \(N = 100\) پاسخ دودویی برازش شده‌اند.
مدل \lr{A:} \(k = 2\)، \lr{NLL} = 60.0. مدل \lr{B:} \(k = 4\)، \lr{NLL}
= 57.5. \lr{AIC} و \lr{BIC} هر دو را حساب کنید. هر معیار کدام را ترجیح
می‌دهد؟ (\(\ln 100 = 4.605\)) \textbf{پاسخ} \lr{AIC: A} = 124.0، \lr{B} =
123.0، پس \lr{AIC} (کمی) \lr{B} را ترجیح می‌دهد. \lr{BIC: A} = 120 + 9.21
= 129.21؛ \lr{B} = 115 + 18.42 = 133.42، پس \lr{BIC A} را ترجیح می‌دهد.
معیارها با هم مخالف‌اند چون \lr{BIC} هر پارامتر را 4.6 به‌جای 2 جریمه
می‌کند.

\textbf{4. (خواندن کد)} این تابع \lr{NLL} چه اشکالی دارد و بهینه‌ساز چه
خواهد کرد؟

\begin{quote}
\textbf{\lr{Code} · کد:} \lr{\nolinkurl{code/week08.py}}, \lr{cell} 4 ·
سلول 4
\end{quote}

\textbf{پاسخ} (1) لگاریتم درست‌نمایی را برمی‌گرداند نه منفی آن را، پس
\lr{\nolinkurl{minimize}} پارامترها را به سمت \emph{بدترین} برازش
(لگاریتم درست‌نمایی به سمت \ensuremath{-}\ensuremath{\infty}) می‌برد. (2) \lr{\nolinkurl{p}} محدود نشده،
پس مقادیر پیش‌بینی‌شده‌ی 0 یا بیشتر از 1 مقدار \lr{\ensuremath{-}inf} یا \lr{nan}
می‌دهند. اصلاح:
\lr{\nolinkurl{return -np.sum(binom.logpmf(y, n, np.clip(p, EPS, 1 - EPS)))}}.

\textbf{5. (تفسیر)} یک مطالعه‌ی بازیابی برای پارامتر \(A\) مقدار
\(r(\text{\lr{true}}, \text{\lr{recovered}}) = 0.95\) و برای پارامتر \(B\) مقدار
\(0.35\) می‌دهد. درباره‌ی مقاله‌ای که با این طرح آزمایش تفاوت گروه‌ها در
\(B\) را گزارش می‌کند چه نتیجه‌ای می‌گیرید؟ \textbf{پاسخ} \(B\) با این طرح
آزمایش ضعیف بازیابی می‌شود: برآوردها تحت تأثیر نویز یا \lr{trade-off}
هستند، پس تفاوت گروه‌ها در \(B\) برآوردشده قابل تفسیر نیست. مقاله باید
طرح آزمایش را بهتر کند یا \(B\) را تفسیر نکند.

\textbf{6. (تفسیر)} یک \lr{confusion matrix BIC} دارای \lr{P(selected} =
\lr{power} \lr{\textbar} \lr{generated} = \lr{asymptote)} = 0.58 است.
برنده‌ی \lr{BIC} در داده‌ی واقعی شما مدل توانی است. چه باید بنویسید؟
\textbf{پاسخ} اینکه \lr{BIC} مدل توانی را ترجیح می‌دهد، ولی با این طرح
آزمایش \lr{BIC} اغلب برای داده‌هایی که مدل مجانبی تولید کرده هم مدل توانی
را انتخاب می‌کند، پس شواهد بین این دو به‌خوبی تمایز نمی‌گذارد. احتمال‌های
معکوس و نتایج \lr{AIC/CV} را گزارش کنید و طرحی با فاصله‌های طولانی‌تر
پیشنهاد دهید.

\textbf{7. (طراحی)} می‌خواهید بدانید آیا یک دارو در مدل \lr{random walk}
احتیاط (مرز \(a\)) را تغییر می‌دهد یا کیفیت شواهد \lr{(drift} \(v\)) را.
چه داده‌ای باید جمع کنید و چرا دقت به‌تنهایی کافی نیست؟ \textbf{پاسخ}
انتخاب‌ها \textbf{و} زمان‌های واکنش (ترجیحاً کل توزیع زمان واکنش) تحت دارو
و دارونما. هم \(a\) بزرگ‌تر و هم \(v\) بزرگ‌تر دقت را بالا می‌برند، ولی
\(a\) پاسخ‌ها را کند و \(v\) سریع می‌کند، پس فقط الگوی مشترک آن‌ها را از هم
جدا می‌کند (هفته‌ی 1).

\textbf{8. (مفهومی)} اگر بخواهیم پیش‌بینی \lr{trial}های جدید را آزمون
کنیم، چرا باید بخش‌های \lr{cross-validation} برای داده‌های فراموشی داخل
فاصله‌ها ساخته شوند؟ \textbf{پاسخ} تا هر مجموعه‌ی آموزش همه‌ی فاصله‌های
نگهداری (شرایط) را داشته باشد و آزمون درباره‌ی آیتم‌های جدید در شرایط
معلوم باشد. کنار گذاشتن کل فاصله‌ها برون‌یابی به شرایط جدید را آزمون می‌کند
که ادعای تعمیم دیگری است.

\begin{center}\rule{0.5\linewidth}{0.5pt}\end{center}

\subsection{تکلیف}

\begin{itemize}
\tightlist
\item
  \textbf{\lr{A3} (مقایسه و بازیابی مدل)} در ابتدای درس هفته‌ی 9 تحویل
  داده می‌شود.
\item
  \textbf{پیشنهاد پروژه} در ابتدای درس هفته‌ی 9 تحویل داده می‌شود (حداکثر
  2 صفحه؛ هفته‌ی 7 را ببینید).
\item
  این هفته تکلیف جدیدی نداریم. برای میان‌ترم با اجرای مطالعه‌ی موردی بالا
  و جواب دادن به سؤال‌های تمرینی بدون نگاه کردن به پاسخ‌ها آماده شوید.
\end{itemize}

\begin{center}\rule{0.5\linewidth}{0.5pt}\end{center}

\subsection{مطالعه‌ی بیشتر}

\begin{itemize}
\tightlist
\item
  \textbf{\lr{Farrell} \lr{\&} \lr{Lewandowsky} (2018)، فصل‌های 2 تا 5 و 10.}
  خلاصه‌ی فصل‌ها را به‌عنوان مرور کل نیمه‌ی اول درس دوباره بخوانید.
\item
  \textbf{\lr{Lee} \lr{\&} \lr{Wagenmakers} (2014)، فصل 11} \lr{(SDT).} توصیف
  مدل را برای آمادگی هفته‌ی 9 بخوانید، که در آن شروع می‌کنیم پارامترها را
  متغیرهای تصادفی با پیشین در نظر بگیریم.
\end{itemize}

\emph{پیش‌نمایش هفته‌ی 9:} نیمه‌ی دوم درس بیزی می‌شود. به‌جای یک بهترین
برآورد به‌علاوه‌ی یک بازه‌ی اطمینان، کل توزیع احتمال پارامترها به شرط داده
را حساب می‌کنیم.
